
\begin{proposition}[The Kuranishi germ of the two-parameter family]\label{cdp:st:E13}
For every member of the family of
\cite[Proposition labelled \texttt{family}]{Viaclovsky2026Two}, on the parameter
bidiscs with $\mu$ sufficiently close to zero and nonzero sufficiently
small linearization scale $\lambda$, one has
\[
 (h^0,h^1,h^2,h^3)(X_{\mu,\lambda},T_X)=(1,2,2,0).
\]
Its Kuranishi germ is a smooth reduced surface, and the displayed
$(\mu,\lambda)$-family is locally miniversal. In particular its
sufficiently small deformations all have algebraic dimension one.
\end{proposition}
\begin{proof}
We work with the equation and section used in the construction,
\[
 S_\mu:\quad y^2=x^3+tx+\mu t+1,\qquad
 P_\mu=(-\mu,b),\qquad b^2=1-\mu^3,
\]
where the branch of $b$ is close to one. Write $p_1,p_2,p_3$ for its
three distinct nodal values, with the zero of the linearization at
$p_3$. The order-four good-reduction cover at infinity gives the
root curve $\mathfrak B=\mathbb P(1,4)$, with coarse projection
$\rho$ and $\rho^*\mathcal O_{\mathbb P^1}(1)=\mathcal O(4)$.
Let $\bar f:X\to\mathfrak B$ denote the root factorization and
$D=W_1+W_2+W_3$ the reduced divisor of Mumford fibres.

We spell out the bundles and the residue map, since their ranks alone
would not determine the answer. The vertical logarithmic Lie bundle
is
\[
 E=\mathcal O\oplus\mathcal O(-1).
\]
The first summand is generated by multiplication in the line-bundle
coordinate. For a concrete splitting of the second summand put
$\omega=dx/(2y)$ and use the relative differential
\begin{equation}\label{cdp:eq:E13-connection}
 \alpha=\frac{y+b}{x+\mu}\,\frac{dx}{2y}.
\end{equation}
Its residues are $1$ at $P_\mu$ and $-1$ at the zero section; there
is no pole at $-P_\mu$, because its numerator vanishes there.
Consequently, in the rational section of
$\mathcal O(P_\mu-O)$ having divisor $P_\mu-O$, the connection
$\nabla(hs)=(dh+\alpha h)s$ is holomorphic along each smooth
elliptic fibre. Here $s$ is that rational section and a point of
the line bundle is written $v s$; the horizontal lift of a vector
$V$ is therefore $\widetilde V=V-\alpha(V)v\partial_v$.
It is a multiplicative connection: the difference between its
pullback under multiplication and the sum of its two pullbacks is a
holomorphic one-form on the product of two elliptic curves, and its
restrictions to both coordinate axes vanish. Every holomorphic
one-form on that product is the sum of forms pulled back from those
axes, so this difference is zero. Its horizontal lift therefore
commutes with the lifted translation defining $X$.

At infinity put $t=r^{-4}$, $x=r^{-2}x'$, $y=r^{-3}y'$. Formula
\eqref{cdp:eq:E13-connection}, in the relative cotangent sheaf, becomes
\[
 \alpha=\frac{y'+br^3}{x'+\mu r^2}\frac{dx'}{2y'},
 \qquad y'^2=x'^3+x'+\mu r^2+r^6.
\]
This is the same holomorphic connection at the good-reduction
fibre; the limiting section $(0,0)$ is a nonzero two-torsion point,
so it remains disjoint from the zero section. Moreover
$\omega=r\,dx'/(2y')$, which gives the summand $\mathcal O(-1)$
for its dual Lie vector. At a nodal fibre the connection is
logarithmic at the two branches and its lift is a section of the
vertical logarithmic tangent bundle. Thus these calculations prove
the claimed splitting on the entire root curve, with the specified
comparison maps.

The inverse pairing associated with the nondegenerate alternating
class $\xi=u\wedge w+2\psi\wedge\delta$ in the construction induces a
holomorphic isomorphism
\[
 G:=R^1\bar f_*\mathcal O_X\simeq E.
\]
Here is the extension check. On a marked smooth fibre the periods
have cross terms in ratio $1:-2$, so the holomorphic one-form
subspace $F^1\subset H^1(F,\mathbb C)$ is Lagrangian for this pairing.
Pairing with $F^1$ therefore gives the isomorphism
$H^1(F,\mathbb C)/F^1\to(F^1)^*=\operatorname{Lie}(F)$.
At a Mumford point, subtract the logarithmic monodromy term from the
periods. In the resulting holomorphic frame the alternating pairing
is still constant and nondegenerate, and the two holomorphic
one-forms are the logarithmic coordinate differentials. Their dual
frame is the logarithmic Lie frame. The quotient frame gives
$R^1\bar f_*\mathcal O_X$: at a rank-one point its first class is
the class of the loop in the normalization diagram and its second
is the class on the elliptic normalization; at the rank-two point
both classes are loops in the normalization diagram. These are
bases, by the normalization exact sequence. Thus the pairing extends
as an isomorphism at these points. The finite good-reduction cover
gives the same conclusion at infinity, and descent uses invariance
of $\xi$. This identifies the comparison, including its value at
the special fibres. Cup product and the normalization sequences
give $R^2\bar f_*\mathcal O_X=\det G=\mathcal O(-1)$.

The logarithmic tangent and canonical sequences are therefore
\[
\begin{split}
 0&\longrightarrow\bar f^*E\longrightarrow T_X(-\log D)
 \longrightarrow\bar f^*\mathcal O(-7)\longrightarrow0,\\
 K_X&=\bar f^*\mathcal O(-4).
\end{split}
\]
Indeed $K_{\mathfrak B}=\mathcal O(-5)$ and
$T_{\mathfrak B}(-p_1-p_2-p_3)=\mathcal O(5-12)$; taking the
determinant of the displayed logarithmic sequence gives the
canonical formula after removing $D$ of degree $12$.
For root line bundles, cohomology is computed by weighted monomials
and Serre duality with $\mathcal O(-5)$. The projection formula and
Leray give
\[
\begin{split}
 h^\bullet(X,\bar f^*E)&=(1,1,0,0),\\
 h^\bullet(X,\bar f^*\mathcal O(-7))&=(0,1,2,1),\\
 h^\bullet(X,T_X(-\log D))&=(1,2,2,1).
\end{split}
\]
There are no uncomputed differentials in these equalities: the base
has dimension one, and the groups of the first displayed row in
degrees two and three vanish.

Let $Q=\bigoplus_{i=1}^3\eta_{i*}K_{\widetilde W_i}$, the normal
quotient of $T_X(-\log D)\to T_X$. Here
$\widetilde W_1,\widetilde W_2$ are ruled over elliptic curves and
$\widetilde W_3=dP_6$, so
\[
 h^\bullet(Q)=(0,2,3).
\]
It remains to compute the boundary
$H^1(Q)\to H^2(T_X(-\log D))$, between two-dimensional spaces.
Under Serre duality its dual is the residue homomorphism
\begin{equation}\label{cdp:eq:E13-residue}
 H^1\bigl(X,\Omega_X^1(\log D)\otimes K_X\bigr)
 \longrightarrow
 H^1(\widetilde W_1,\mathcal O)\oplus
 H^1(\widetilde W_2,\mathcal O).
\end{equation}
The dual logarithmic sequence, tensored with $K_X$, has horizontal
term $\bar f^*\mathcal O(3)$ and vertical quotient
$\bar f^*(\mathcal O(-4)\oplus\mathcal O(-3))$.
The latter has zero cohomology in degrees zero and one: the coarse
direct images of $\mathcal O(-4)$ and $\mathcal O(-3)$ are both
$\mathcal O_{\mathbb P^1}(-1)$, and every line in their tensor
products with $G$ has negative weighted degree, so has no section.
Hence the
source of \eqref{cdp:eq:E13-residue} is canonically
\[
 H^1(X,\bar f^*\mathcal O(3))
 =H^0\bigl(\mathfrak B,\mathcal O(3)\otimes G\bigr)
 =H^0(\mathcal O(3))\oplus H^0(\mathcal O(2))=\mathbb C^2.
\]
On the affine $t$-line the generators of these last two spaces are
constant in the frame just constructed. The residue in each
summand is a nonzero scalar times evaluation followed by
$G_{p_i}\to H^1(\widetilde W_i,\mathcal O)$.

Write $e_i$ for the double root in the nodal cubic at $p_i$.
The line killed by this last projection is the graph-loop class.
Under $G\simeq E$ it is the toric Lie line: in the local periods
$(q\alpha,\beta),(\gamma,\lambda\delta)$ it is the direction
of the degenerating first multiplicative coordinate, whereas the
second coordinate is the elliptic normalization. If $V$ is dual to
$\omega$, its horizontal lift $\widetilde V$ has vertical
line-coordinate component $-\alpha(V)$. At the nodal boundary put
\[
 a_i=\frac{b}{e_i+\mu}.
\]
Thus the toric line is generated by
$\widetilde V+a_i v\partial_v$, and the projection to the elliptic
Lie line sends $(v\partial_v,\widetilde V)$ to $(1,-a_i)$, up to
a common nonzero factor at that point. Consequently the matrix of
\eqref{cdp:eq:E13-residue}, after nonzero row rescalings, is
\[
 \begin{pmatrix}1&-a_1\\1&-a_2\end{pmatrix},\qquad
 a_1-a_2=\frac{b(e_2-e_1)}{(e_1+\mu)(e_2+\mu)}\ne0.
\]
The inequality is exact: $b\ne0$; distinct nodal values have
distinct double roots, since $p_i=-3e_i^2$; and $e_i=-\mu$ would
place $P_\mu$ at the node and force $b=0$. The boundary therefore
has rank two.

The rank-two cusp implies
$\bar f_*\Omega_X^1=\Omega_{\mathfrak B}^1$, by the generic
Kodaira--Spencer boundary and extension across the local models,
as in \ref{cdp:st:E11}. It follows that
$h^3(T_X)=h^0(\Omega_X^1\otimes K_X)
=h^0(\mathcal O(-9))=0$. The long exact sequence with quotient
$Q$ now gives $h^0(T_X)=1$, $h^1(T_X)=2$, and $h^2(T_X)=2$.

Finally the constructed bidisc has injective Kodaira--Spencer map.
At $\mu=0$, writing $a_i=p_i(0)$, differentiation of
$4t^3+27(\mu t+1)^2=0$ gives $p_i'(0)=-9/(2a_i)$.
An infinitesimal automorphism fixing the multiple value at infinity
has the form $At+B$; the three equalities
$-9/(2a_i)=Aa_i+B$ would make a nonzero quadratic polynomial vanish
at all three $a_i$. The cross-ratio variation is therefore nonzero,
and remains so on a smaller disc. For fixed $\mu$, a normalized
double curve is $\mathbb C^*/\langle c(\mu)\lambda\rangle$.
Its elliptic modulus has nonzero $\lambda$-derivative when
$|c(\mu)\lambda|$ is sufficiently small, as follows from
$j(q)=q^{-1}+744+O(q)$. A trivial first-order deformation of $X$
would trivialize its intrinsic anticanonical pencil, its labelled
exceptional-value configuration, and this double curve; it can
therefore kill neither of the two independent tangent directions.
Thus the bidisc maps to the two-dimensional tangent space of the
Kuranishi germ with invertible differential. Present its local
ring as $\mathbb C\{z_1,z_2\}/I$; the induced substitution into
$\mathbb C\{\mu,\lambda-\lambda_0\}$ is a local analytic
isomorphism before quotienting, by the inverse function theorem.
Every element of $I$ vanishes under it, so $I=0$. This proves the
Kuranishi assertion. \ref{cdp:st:E1} and \ref{cdp:st:E6} give the algebraic-dimension assertion.
\end{proof}


